\mychapter{Objetivos}{cap:objetivos}
\lhead{OBJETIVOS}

\textbf{Objetivo geral.} Avaliar se aplicações geradas a partir de uma fonte única, com projetos nativos por sistema operacional, atingem desempenho comparável ao de aplicações escritas à mão nas ferramentas nativas de Windows e macOS.

\textbf{Objetivos específicos.}
\begin{enumerate}
  \item Quantificar a diferença de desempenho (tempo de resposta, memória e CPU) entre as aplicações geradas e as escritas à mão.
  \item Comparar as aplicações geradas com uma versão de renderização própria.
  \item Medir a similaridade visual das aplicações geradas em relação às escritas à mão.
  \item Determinar a fração dos componentes de interface de cada sistema operacional que a fonte única consegue expressar.
\end{enumerate}

\textbf{Pergunta de pesquisa.} Como desenvolver aplicativos multiplataforma que fazem uso dos componentes nativos de interface em uma base única de código, sem perda de desempenho em relação ao desenvolvimento nativo?

\textbf{Hipóteses.} Formuladas antes dos experimentos, H1 a H3 competem entre si.
\begin{itemize}
  \item \textbf{H1:} as aplicações geradas têm desempenho equivalente ao das escritas à mão, isto é, a diferença entre as medianas fica dentro de uma margem de equivalência pré-definida.
  \item \textbf{H2:} as aplicações geradas têm desempenho superior ao da versão com renderização própria, tanto no Windows quanto no macOS.
  \item \textbf{H3 (rival):} o custo da geração e da camada compartilhada anula o ganho, e as aplicações geradas ficam equivalentes à versão com renderização própria.
  \item \textbf{H4 (cobertura):} a fonte única não cobre todos os componentes nativos do Windows e macOS, e a similaridade visual (SSIM) é menor nas telas com componentes não cobertos.
\end{itemize}
